\documentclass[12pt]{article}

\usepackage[a4paper,margin=0.9in]{geometry}
\usepackage{amsmath,amssymb}
\usepackage{newtxtext,newtxmath} % Times-style font
\usepackage{fancyhdr}
\usepackage{listings}
\usepackage{graphicx}
\usepackage{float}
\usepackage{booktabs}
\usepackage{hyperref}

\setlength{\parindent}{0pt}
\setlength{\parskip}{5pt}

% Keep entire document at 12 pt
\AtBeginDocument{
    \fontsize{12pt}{14pt}\selectfont
}

% ------------------------------------------------
% Header
% ------------------------------------------------

\pagestyle{fancy}
\fancyhf{}

\fancyhead[L]{Name: \underline{Shaik Jnaved}}
\fancyhead[R]{Enrollment No.: \underline{P26DS001}}

\renewcommand{\headrulewidth}{0.4pt}
\setlength{\headheight}{15pt}

% ------------------------------------------------
% Code formatting
% ------------------------------------------------

\lstset{
    basicstyle=\ttfamily\fontsize{9pt}{11pt}\selectfont,
    breaklines=true,
    frame=single,
    columns=fullflexible,
    keepspaces=true,
    showstringspaces=false
}

\begin{document}

% ------------------------------------------------
% TITLE
% ------------------------------------------------

\begin{center}
\textbf{Computer Science and Engineering Department, SVNIT, Surat}\\
\textbf{M.Tech. I -- Semester I}\\
\textbf{Computer Vision and Image Processing (CS6111 / CSDS119)}\\
\textbf{Tool Exploration Report: MeshLab and Autodesk Meshmixer}
\end{center}

\vspace{0.2cm}

\textbf{Problem Statement:} \\
Explore the assigned tools (\textbf{Tool 1: MeshLab} and \textbf{Tool 2: Meshmixer}) for 3D Computer Vision and Image Processing. Perform a small, meaningful experiment on each tool, demonstrate input/output results, formulate the underlying CV/IP theoretical concepts, and document critical observations and limitations.

\vspace{0.3cm}
\hrule
\vspace{0.3cm}

\section*{1. Tool 1: MeshLab}

\subsection*{1.1 Purpose and CV/IP Application}
MeshLab is an open-source, portable, and extensible 3D processing system designed for editing and processing unstructured 3D triangular meshes and dense 3D point clouds. In Computer Vision, its primary applications include \textbf{3D Surface Reconstruction from unorganized LiDAR/RGB-D point clouds}, spatial normal estimation, and \textbf{Geometric Mesh Simplification (Decimation)}.

\subsection*{1.2 Underlying CV/IP Concepts}
\begin{enumerate}
    \item \textbf{Screened Poisson Surface Reconstruction (Kazhdan et al.):} Formulates surface reconstruction as a spatial Poisson boundary-value problem:
    \begin{equation}
        \Delta \chi = \nabla \cdot \vec{V}
    \end{equation}
    where $\chi : \mathbb{R}^3 \to \mathbb{R}$ is an indicator function ($1$ inside the object, $0$ outside), and $\vec{V}$ is the smoothed vector field of surface normals. Solving this on an adaptive octree recovers an implicit surface extracted using Marching Cubes.
    \item \textbf{Quadric Error Metric (QEM) Decimation (Garland \& Heckbert):} Iteratively collapses edges $(v_1, v_2) \to \bar{v}$ by minimizing the quadratic sum of squared distances to adjacent triangle planes:
    \begin{equation}
        \Delta(\bar{v}) = \bar{v}^T (Q_1 + Q_2) \bar{v}
    \end{equation}
\end{enumerate}

\subsection*{1.3 Practical Experiment \& Results}
\begin{itemize}
    \item \textbf{Input:} A dense unorganized 3D point cloud of $12,000$ points with synthetic Gaussian acquisition noise ($\sigma = 0.006$).
    \item \textbf{Step 1 (Poisson Reconstruction):} Reconstructed a watertight, 2-manifold surface comprising $72,924$ triangular faces.
    \item \textbf{Step 2 (QEM Decimation):} Simplified geometry to $3,500$ faces (a $95.2\%$ reduction in complexity) while preserving high-curvature silhouettes.
\end{itemize}

\begin{figure}[H]
    \centering
    \includegraphics[width=\textwidth]{images/meshlab_experiment_comparison.png}
    \caption{MeshLab Experiment: (A) Input Point Cloud with Noise, (B) Screened Poisson Surface, (C) QEM Decimated Mesh with Wireframe.}
\end{figure}

\textbf{Key Observation / Limitation:} Poisson surface reconstruction is highly sensitive to the consistency and orientation of normal vectors. Flipped or noisy normals produce artificial bubbling or self-intersecting topological defects.

\newpage

\section*{2. Tool 2: Autodesk Meshmixer}

\subsection*{2.1 Purpose and CV/IP Application}
Autodesk Meshmixer is a specialized tool for 3D mesh modeling, inspection, and sculpting. In 3D Vision, it is widely utilized for \textbf{3D Scan Defect Inspection, Topological Hole Closure, and Isotropic Remeshing} to prepare raw 3D scans (e.g., medical CT/MRI reconstructions, photogrammetry outputs) for downstream numerical analysis and fabrication.

\subsection*{2.2 Underlying CV/IP Concepts}
\begin{enumerate}
    \item \textbf{2-Manifold Geometry \& Euler Characteristic:} A 2-manifold mesh requires every edge to be shared by at most two triangles. By Euler's Polyhedral Formula, a closed surface of genus $g$ satisfies:
    \begin{equation}
        \chi = V - E + F = 2 - 2g \quad (\chi = 2 \text{ for a watertight closed sphere})
    \end{equation}
    \item \textbf{Topological Boundary Loop Healing:} Open scan boundaries (edges with face count $= 1$) are detected through the 1-ring boundary graph and filled via minimal-energy Delaunay patch triangulation.
    \item \textbf{Isotropic Remeshing \& Discrete Laplacian Fairing:} Smoothing moves vertices along their discrete Laplace-Beltrami operator to eliminate high-frequency scanning noise, while remeshing equalizes triangle edge lengths to eliminate ill-conditioned sliver elements.
\end{enumerate}

\subsection*{2.3 Practical Experiment \& Results}
\begin{itemize}
    \item \textbf{Input:} Defective 3D scan with $6$ open boundary holes, $176$ boundary edges, and non-watertight topology ($\chi = -4$).
    \item \textbf{Step 1 (Inspection \& Hole Closure):} Automatically detected boundary cycles, constructed planar/curved Delaunay patches, and restored the mesh to a watertight 2-manifold state ($\chi = 2$, $5,252$ faces).
    \item \textbf{Step 2 (Isotropic Remeshing):} Regularized triangulation to $5,000$ uniform equilateral triangles.
\end{itemize}

\begin{figure}[H]
    \centering
    \includegraphics[width=\textwidth]{images/meshmixer_experiment_comparison.png}
    \caption{Meshmixer Experiment: (A) Defective 3D Scan with Open Holes, (B) Watertight Repaired Mesh ($\chi=2$), (C) Isotropic Remeshed \& Faired Surface.}
\end{figure}

\textbf{Key Observation / Limitation:} Meshmixer excels at interactive manual sculpting and repair brushes, but lacks programmable headless/batch scripting interfaces compared to MeshLab or Open3D.

\vspace{0.3cm}
\hrule
\vspace{0.3cm}

\section*{3. Quantitative Comparison \& Summary}

\begin{table}[H]
\centering
\begin{tabular}{@{}lccc@{}}
\toprule
\textbf{Pipeline Stage} & \textbf{Face Count} & \textbf{Watertight Status} & \textbf{Euler Characteristic ($\chi$)} \\ \midrule
\textit{MeshLab: Point Cloud (Input)} & $12,000$ pts & N/A & N/A \\
\textit{MeshLab: Poisson Reconstructed} & $72,924$ & \textbf{True} & $2$ \\
\textit{MeshLab: QEM Decimated ($95.2\%$)} & $3,500$ & \textbf{True} & $2$ \\ \midrule
\textit{Meshmixer: Defective Scan (Input)} & $5,076$ & \textbf{False} & $-4$ \\
\textit{Meshmixer: Hole-Filled \& Healed} & $5,252$ & \textbf{True} & \textbf{2} \\
\textit{Meshmixer: Isotropic Remeshed} & $5,000$ & \textbf{True} & $2$ \\ \bottomrule
\end{tabular}
\caption{Quantitative summary of geometric and topological metrics across both experiments.}
\end{table}

\begin{figure}[H]
    \centering
    \includegraphics[width=0.88\textwidth]{images/quantitative_comparison_dashboard.png}
    \caption{Comparative analysis chart of mesh complexity across both experimental pipelines.}
\end{figure}

\end{document}
